(1) 无损细长号角管道截面积为 \(A(x)=A_0 e^{2\alpha x}\)，\(\alpha>0\)，空气密度 \(\rho\) 和声速 \(c\) 均为常数，波长长于横向尺度。由连续性和一维无黏动量方程推导压力扰动 \(p(x,t)\) 的Webster方程，再用 \(p=e^{-\alpha x}u\) 求等效定系数波动方程。
(2) 对无限指数号角，令 \(u=e^{i(kx-\omega t)}\) 推导频率、截止条件、相速及群速。说明截止以下的压力场如何倏逝，及传播波的压力振幅为什么随 \(x\) 变。
(3) 从基本声学方程导出局部声能守恒。对单一向右传播模计算周期平均压力、质点速度、声能密度及轴向能流，证明平均能流/能量密度等于群速度，且总通过功率 \(A(x)\overline{pv}\) 不随位置变。
(4) 有限区间 \(0<x<L\) 两端均是刚壁，即 \(v(0,t)=v(L,t)=0\)。求非零固有频率与压力本征函数；区分物理的零频均匀压力模，并说明正交权函数及不同本征频率振型的正交关系。
